% Bagian: first-order-logic
% Bab: models-theories
% Subbagian: introduction

\documentclass[../../../include/open-logic-section]{subfiles}

\begin{document}

\olfileid[id]{fol}{mat}{int}
\olsection{Pendahuluan}

\begin{explain}
Pengembangan metode aksiomatik merupakan pencapaian penting dalam sejarah
ilmu pengetahuan, dan sangat penting khususnya dalam sejarah matematika.
Pengembangan aksiomatik suatu bidang melibatkan penjernihan banyak pertanyaan:
Apa yang dibahas oleh bidang itu? Apa konsep-konsepnya yang paling mendasar?
Bagaimana konsep-konsep tersebut saling berkaitan? Dapatkah semua konsep dalam
bidang itu didefinisikan berdasarkan konsep-konsep mendasar tersebut? Hukum apa
yang dipatuhi, dan harus dipatuhi, oleh konsep-konsep ini?

Metode aksiomatik dan logika sangat serasi. Logika formal menyediakan
perangkat untuk merumuskan teori aksiomatik, untuk membuktikan teorema dari
aksioma-aksioma teori dengan cara yang ditentukan secara saksama, dan untuk
mengkaji secara sistematis sifat-sifat semua sistem yang memenuhi aksioma
tersebut.
\end{explain}

\begin{defn}
Suatu himpunan !!{sentence}s~$\Gamma$ disebut \emph{tertutup} jika dan hanya
jika, setiap kali $\Gamma \Entails !A$, berlaku $!A \in \Gamma$.
\emph{Penutupan} suatu himpunan !!{sentence}s~$\Gamma$ ialah
$\Setabs{!A}{\Gamma \Entails !A}$.

Kita mengatakan bahwa~$\Gamma$ \emph{diaksiomatisasi oleh} suatu himpunan
kalimat~$\Delta$ jika $\Gamma$ merupakan penutupan dari~$\Delta$.
\end{defn}

\begin{explain}
Kita dapat memandang suatu teori aksiomatik sebagai himpunan kalimat yang
diaksiomatisasi oleh himpunan aksiomanya~$\Delta$. Dengan kata lain, apabila
kita mempunyai bahasa orde pertama yang memuat simbol nonlogis bagi
konsep-konsep primitif dari ilmu yang dikembangkan secara aksiomatik yang ingin
kita kaji, beserta suatu himpunan !!{sentence}s yang mengungkapkan hukum-hukum
mendasar ilmu tersebut, kita dapat memandang teori itu sebagai direpresentasikan
oleh semua !!{sentence}s dalam bahasa ini yang diakibatkan oleh aksioma-aksioma
tersebut. Contohnya berkisar dari teori sederhana dengan hanya satu konsep
primitif dan aksioma sederhana, seperti teori urutan parsial, hingga teori
rumit seperti mekanika Newton.

Fakta-fakta logis penting yang membuat pendekatan formal terhadap metode
aksiomatik ini begitu penting adalah sebagai berikut. Andaikan $\Gamma$ suatu
sistem aksioma bagi sebuah teori, yaitu suatu himpunan kalimat.
\begin{enumerate}
\item Kita dapat menyatakan secara saksama kapan suatu sistem aksioma menangkap
  kelas !!{structure}s yang dimaksudkan. Yakni, jika kita tertarik pada suatu
  kelas !!{structure}s tertentu, kita berhasil menangkap kelas tersebut dengan
  sistem aksioma~$\Gamma$ jika dan hanya jika !!{structure}s itu tepat
  semua~$\Struct M$ sedemikian sehingga $\Sat{M}{\Gamma}$.
\item Kita mungkin gagal dalam hal ini karena terdapat $\Struct M$ sedemikian
  sehingga $\Sat{M}{\Gamma}$, tetapi $\Struct M$ bukan salah satu
  !!{structure}s yang kita maksudkan. Hal ini dapat mendorong kita untuk
  menambahkan aksioma yang tidak benar dalam~$\Struct M$.
\item Jika kita berhasil setidaknya dalam arti bahwa $\Gamma$ benar dalam
  semua !!{structure}s yang dimaksudkan, maka suatu kalimat~$!A$ benar dalam
  semua !!{structure}s yang dimaksudkan setiap kali $\Gamma \Entails !A$.
  Dengan demikian, kita dapat menggunakan perangkat logis (seperti metode
  !!{derivation}) untuk menunjukkan bahwa kalimat-kalimat benar dalam semua
  !!{structure}s yang dimaksudkan, cukup dengan menunjukkan bahwa kalimat itu
  diakibatkan oleh aksioma-aksioma tersebut.
\item Kadang-kadang kita tidak mempunyai !!{structure}s yang dimaksudkan dalam
  pikiran, tetapi justru berangkat dari aksioma itu sendiri: kita mulai dengan
  beberapa konsep primitif yang kita kehendaki mematuhi hukum tertentu, yang
  kita kodifikasikan dalam suatu sistem aksioma. Salah satu hal yang ingin kita
  verifikasi segera ialah bahwa aksioma-aksioma itu tidak saling bertentangan:
  jika saling bertentangan, tidak mungkin ada konsep yang mematuhi hukum-hukum
  tersebut, dan kita telah mencoba membangun teori yang tidak koheren. Kita
  dapat memverifikasi bahwa hal ini tidak terjadi dengan menemukan suatu model
  dari~$\Gamma$. Dan jika teori kita mempunyai model, kita dapat menggunakan
  metode logis untuk menyelidikinya, serta untuk membangun model.
\item Independensi aksioma juga merupakan pertanyaan penting. Salah satu
  aksioma mungkin sebenarnya merupakan konsekuensi dari aksioma-aksioma lain,
  sehingga aksioma itu berlebihan. Kita dapat membuktikan bahwa suatu aksioma
  $!A$ dalam $\Gamma$ berlebihan dengan membuktikan $\Gamma \setminus
  \{!A\} \Entails !A$. Kita juga dapat membuktikan bahwa suatu aksioma tidak
  berlebihan dengan menunjukkan bahwa $(\Gamma \setminus \{!A\}) \cup
  \{\lnot !A\}$ dapat dipenuhi. Misalnya, dengan cara inilah postulat sejajar
  ditunjukkan independen dari aksioma-aksioma geometri lainnya.
\item Pertanyaan penting lainnya ialah keterdefinisian konsep dalam suatu
  teori: Pilihan bahasa menentukan apa saja yang membentuk model-model suatu
  teori. Namun, tidak setiap aspek teori harus direpresentasikan secara
  terpisah dalam modelnya. Misalnya, setiap urutan $\le$ menentukan urutan
  ketat~$<$ yang bersesuaian---jika salah satunya diberikan, kita dapat
  mendefinisikan yang lain. Jadi, model suatu teori yang melibatkan urutan
  semacam itu tidak harus \emph{juga} memuat urutan ketat yang bersesuaian.
  Secara umum, kapan suatu relasi dapat didefinisikan berdasarkan relasi-relasi
  lain? Kapan mustahil mendefinisikan suatu relasi berdasarkan relasi-relasi
  lain (sehingga relasi itu harus ditambahkan pada konsep-konsep primitif
  bahasa)?
\end{enumerate}
\end{explain}


\end{document}
